\documentclass[10pt,letterpaper]{article}

\usepackage{cmap}
\usepackage[margin=0.68in]{geometry}
\usepackage[T1]{fontenc}
\usepackage[utf8]{inputenc}
\usepackage{array}
\usepackage{enumitem}
\usepackage[scaled=0.96]{helvet}
\usepackage{hyperref}
\usepackage{microtype}
\usepackage{ragged2e}
\usepackage{tabularx}
\usepackage{xcolor}
\usepackage{titlesec}

\definecolor{accent}{HTML}{B86F63}
\definecolor{accenttwo}{HTML}{8F524A}
\definecolor{text}{HTML}{261712}
\definecolor{muted}{HTML}{715D56}
\definecolor{peach}{HTML}{FFF0EB}
\definecolor{rule}{HTML}{DABEB7}

\hypersetup{
  colorlinks=true,
  urlcolor=accent,
  linkcolor=accent
}

\renewcommand{\familydefault}{\sfdefault}
\pagestyle{empty}
\raggedbottom
\setlength{\parindent}{0pt}
\setlength{\parskip}{0pt}
\hyphenpenalty=10000
\exhyphenpenalty=10000
\clubpenalty=10000
\widowpenalty=10000
\displaywidowpenalty=10000
\brokenpenalty=10000
\setlist[itemize]{label=\textcolor{accent}{\small\textbullet},leftmargin=1.25em,itemsep=2.6pt,topsep=3pt,parsep=0pt}
\titleformat{\section}{\large\bfseries\color{accent}}{}{0pt}{}[\vspace{-2pt}{\color{rule}\titlerule}]
\titlespacing{\section}{0pt}{5pt}{2.5pt}

\newcommand{\contactsep}{\quad\textcolor{muted}{\textbullet}\quad}
\newcommand{\cvtag}[1]{\begingroup\setlength{\fboxsep}{3.2pt}\colorbox{peach}{\textcolor{accent}{\footnotesize\bfseries #1}}\endgroup\hspace{3pt}}
\newcommand{\profilebox}[1]{%
  \vspace{5pt}
  {\setlength{\fboxsep}{8pt}\colorbox{peach}{\parbox{\dimexpr\linewidth-2\fboxsep\relax}{#1}}}
  \vspace{1pt}
}
\newcommand{\roleentry}[4]{%
  \begin{tabularx}{\textwidth}{@{}>{\raggedright\arraybackslash}X r@{}}
    \textbf{#1} & \textbf{#2}\\
    \textit{\textcolor{muted}{#3}} & \textit{\textcolor{muted}{#4}}
  \end{tabularx}
  \vspace{-2pt}
}
\newcommand{\degreeentry}[4]{%
  \begin{tabularx}{\textwidth}{@{}>{\raggedright\arraybackslash}X r@{}}
    \textbf{#1} & \textbf{#2}\\
    #3 & #4
  \end{tabularx}
  \vspace{0.25pt}
}
\newcommand{\datedline}[2]{%
  \begin{tabularx}{\textwidth}{@{}>{\raggedright\arraybackslash}X r@{}}
    #1 & #2
  \end{tabularx}
}

\begin{document}
\RaggedRight

{\fontsize{22}{25}\selectfont\textbf{Alan Nguyen Pham}}\\[-1pt]
{\large\textbf{\textcolor{accent}{Soft Robotics R\&D \textbar{} Morphing Surfaces \textbar{} Computational Design}}}\\[-1pt]
\href{tel:7142274396}{(714) 227-4396}
\contactsep \href{mailto:anpham@wpi.edu}{anpham@wpi.edu}
\contactsep Worcester, MA
\contactsep \href{https://aolabs.io}{aolabs.io}
\contactsep \href{https://www.linkedin.com/in/alanpham97/}{linkedin.com/in/alanpham97}\\[4pt]
\profilebox{\textcolor{text}{Mechanical engineering Ph.D. candidate developing soft robotic morphing surfaces that transform flat sheets into 3D forms with overhangs and wave-like profiles. Work combines computational design, compliant mechanisms, pneumatic actuation, modular surface architectures, and physical prototyping to turn target shapes into buildable hardware.}}

\section*{Research Focus}
\textcolor{muted}{Soft robotic morphing surfaces; overhang formation; wave-like geometry; computational design; compliant mechanism design; pneumatic actuation and routing; modular surface architectures; snap-fit pneumatic-mechanical connections; simulation-guided prototyping; physical hardware iteration.}

\section*{Education}
\degreeentry{Doctor of Philosophy, Mechanical Engineering}{Expected 2027}{Worcester Polytechnic Institute, Worcester, MA}{Advisors: Dr. Pratap Rao and Dr. Cagdas Onal}
\degreeentry{Master of Science, Mechanical Engineering}{2022}{California State University, Long Beach, Long Beach, CA}{Advisor: Dr. Allen Teagle-Hernandez}
\degreeentry{Bachelor of Science, Mechanical Engineering}{2019}{California State University, Maritime Academy, Vallejo, CA}{Advisor: Dr. Tomas Oppenheim}

\section*{Awards}
\datedline{NSF NRT Future of Robots in the Workplace -- Research and Development (FORW-RD) Fellowship}{2024}
\datedline{WPI Intercollegiate Soft Robotics Competition -- 1st Place}{2024}

\section*{Research Experience}
\roleentry{Worcester Polytechnic Institute, Worcester, MA}{Aug 2023 -- Present}{Graduate Researcher -- NanoEnergy Lab and Soft Robotics Lab}{Advisors: Dr. Pratap Rao and Dr. Cagdas Onal}
\begin{itemize}
  \item Developing soft robotic morphing surfaces that turn flat cell arrays into overhangs, waves, domes, and moving 3D contours by tuning cell geometry, displacement, and connector layout.
  \item Built monolithically printed Sarrus/PneuNet-inspired cells with cross-cell air channels, bending bias, snap-fit connectors, and planar-to-cylinder assembly paths.
  \item Demonstrated single/double domes, object transport/capture, peristaltic crawling, and self-rolling locomotion from the same cell architecture.
  \item Created Wavevis, a computational design tool that maps target overhang/wave shapes to double-layer Sarrus lattice topology, cell displacements, connector geometry, and prototype dimensions.
  \item Developing FluxCell around one shared air inlet, self-sealing snap-fit pneumatic connections, and magnet/EPM valves that route manifold pressure into printed PneuNets.
\end{itemize}

\roleentry{NASA L' SPACE (NPWEE), Tempe, AZ}{Aug 2025 -- Present}{Soft Robotics Graduate Researcher -- Project SELENE}{}
\begin{itemize}
  \item Contributed to Project SELENE, a NASA NPWEE effort developing soft robotic arms for lunar terrain traversal and regolith sample collection.
  \item Developed a SpiRobs-inspired modular arm with logarithmic-spiral segments, triple-cable actuation, and an elastic braided-steel backbone for curling and grasping motions.
  \item Iterated CAD from two-cable unit concepts to a three-cable spiral architecture, sizing cable/core passages, support strategy, and motor-spool packaging around printable segments.
  \item Prototyped 25\% and full-scale PLA/PETG arm segments, fit-tested actuation cables and elastic cores, and adjusted hole sizes, tolerances, and support placement for assembly.
  \item Produced a full-scale ten-segment PETG arm with inserted cables; evaluated hand actuation, cable rigidity, spool fit, and motorized-test constraints while coordinating with teammates and NASA subject matter experts.
\end{itemize}

\roleentry{Sandia National Laboratories, Albuquerque, NM}{June 2025 -- Aug 2025}{Graduate Research Intern -- Nonlinear Mechanics and Dynamics (NOMAD) Research Institute}{}
\begin{itemize}
  \item Co-developed NOMAD finite-element sweeps for an electromechanical ratcheting mechanism with intermittent contact, friction, springs, and pawl-gear interactions.
  \item Built comparable analytical, MATLAB multibody, and finite-element submodels for pin-spring-pawl, pin-pawl, and pawl-gear interactions under haversine shock and sinusoidal vibration.
  \item Ran structured sweeps across momentum-balance iterations (1--100), mesh density, and processor count (50--250\%) to separate physical mechanism response from solver and modeling sensitivity.
  \item Identified sensitivity regimes in which pin-spring-pawl models converged while pin-pawl and pawl-gear cases showed contact/friction divergence, mesh-related lock-up, and environment-dependent sensitivity.
  \item Presented final results showing submodel-scale simulation sensitivity, informing mesh, contact/friction, iteration, and compute-resource choices before full assembly-level ratcheting models.
\end{itemize}

\roleentry{Worcester Polytechnic Institute, Worcester, MA}{Aug 2023 -- Present}{NSF NRT FORW-RD Fellow and Trainee}{}
\begin{itemize}
  \item Selected as an NSF NRT FORW-RD trainee in an interdisciplinary program on robotic interfaces and assistants for future workplaces spanning robotics, mechanical engineering, computer science, materials, business, humanities, and social sciences.
  \item Connected doctoral research to virtual training, soft robotic teleoperation, compliant interaction, and robotic systems for high-skill physical tasks.
  \item Wrote broader-impact material for robotics and AI systems that operate near people, emphasizing workplace adoption, trust, training, and responsible deployment.
  \item Secured a \$1,000 grant supporting equipment, travel, and professional-development costs.
\end{itemize}

\roleentry{Worcester Polytechnic Institute, Worcester, MA}{Aug 2023 -- Dec 2023}{Research Assistant -- Cognitive Medical Technology Laboratory}{}
\begin{itemize}
  \item Fabricated and assembled low-cost open-source actuation hardware for concentric tube continuum robots, a miniaturized robot class for precise navigation in constrained spaces and minimally invasive surgery.
  \item Contributed to accessible CTR hardware with easy-to-source components, bills of materials, assembly instructions, and guided learning materials for early graduate students.
  \item Integrated actuation hardware, MATLAB forward kinematics, and optical-tracking experiments to validate end-effector displacement.
  \item Coordinated sourcing, fabrication, and assembly across WPI's Cognitive Medical Technology Lab, PracticePoint, and Innovation Studio for two complete CTR systems.
\end{itemize}

\roleentry{California State University, Long Beach, Long Beach, CA}{Sept 2022 -- Aug 2023}{Research Assistant -- Multiscale Mechanics of Materials Group}{Advisor: Dr. Mortaza Saeidi-Javash}
\begin{itemize}
  \item Built a literature foundation for a flexible-electronics research program covering printed and conformal sensors, 2D nanomaterial inks, and wearable human-machine-interface devices.
  \item Synthesized fabrication approaches across inkjet, screen, aerosol-jet, extrusion, and direct-ink-write printing using graphene, MXene, MoS2, thermoelectric, and piezoelectric material systems.
  \item Evaluated sensitivity, selectivity, response time, mechanical compliance, and durability of flexible physical, chemical, and biosensors for soft robotics, structural health monitoring, and wearable systems.
  \item Collaborated with graduate researchers on a direct-ink-write temperature sensor printed on a flexible PET substrate.
\end{itemize}

\roleentry{California State University, Long Beach, Long Beach, CA}{April 2021 -- July 2022}{Graduate Researcher -- Violin Acoustics and Structural Dynamics}{Advisor: Dr. Allen Teagle-Hernandez}
\begin{itemize}
  \item Completed a master's thesis translating perceived violin quality into measurable vibroacoustic features, linking structural design, harmonic balance, frequency response, and performer/listener perception.
  \item Designed a controlled comparison of high- and low-cost modern violins, standardizing strings and fittings while documenting top plate, bridge, f-hole, bass bar, sound post, and tailpiece variables.
  \item Built a LabVIEW/NI data-acquisition workflow with PCB accelerometers, an array microphone, and an impulse modal hammer; collected 67-point top-plate maps and anechoic-chamber acoustic measurements.
  \item Linked FFT-derived violin response features to modal balance and perceived tone.
\end{itemize}

\roleentry{California State University, Maritime Academy, Vallejo, CA}{Sept 2018 -- May 2019}{Undergraduate Researcher -- Cross-Flow Tidal Turbine}{Advisor: Dr. Tomas Oppenheim}
\begin{itemize}
  \item Conceptualized a Darrieus cross-flow tidal turbine, defined rotor geometry and build constraints, and selected 3D-printed blades with fiberglass coating as the manufacturing approach.
  \item Verified performance through water-tunnel power assessment and ANSYS structural analysis, evaluating blade and hub stresses under hydrodynamic loading and documenting operation around 25 rpm.
  \item Fabricated a full-scale prototype using additive manufacturing and sheet-metal forming informed by scaled water-tunnel experiments.
\end{itemize}

\section*{Publications and Research Outputs}
\begin{itemize}
  \item Pham, A. N., Rao, P., \& Onal, C. D. (2026). \textit{A monolithic cellular architecture for topology-programmable soft robotics}. Manuscript in preparation.
  \item Winston, S., Pham, A., Bahr, B., Singh, A., Schumann, C., Flicek, R., Grutzik, S., Kuether, R. J., \& Vemaganti, K. (2025). \href{https://aolabs.io/sandia/paper.pdf}{\textit{Quantifying the Effect of Non-Physical Parameters on the Nonlinear Dynamics of an Electromechanical Ratcheting Mechanism}}. Accepted to IMAC, A Conference and Exposition on Structural Dynamics.
  \item Bonofiglio, K. (2023). \href{https://digital.wpi.edu/pdfviewer/ms35td790}{\textit{Making Concentric Tube Robots More Accessible: An Open-Source Platform with Easy-to-Source Components and Guided Learning Materials}} [Master's thesis, Worcester Polytechnic Institute]. Research contributor to open-source CTR hardware assembly and validation support.
  \item Pham, A., \& Teagle Hernandez, A. (2022). \href{https://www.proquest.com/docview/2718992967}{\textit{Structural Properties and Perceived Acoustic Quality of Modern Violins}} (Publication No. 29065875) [Master's thesis, California State University, Long Beach]. ProQuest Dissertations and Theses Global.
\end{itemize}

\section*{Professional Experience}
\roleentry{Worcester Polytechnic Institute, Worcester, MA}{Aug 2023 -- Present}{Teaching Assistant}{}
\begin{itemize}
  \item Developed learning activities for Stress Analysis, Statics, and Computer Aided Design with professors and teaching teams.
  \item Led weekly conference sessions and office hours for 90+ student groups, translating difficult engineering material into interactive problem-solving sessions.
  \item Designed hands-on learning sessions and guided students in connecting classroom mechanics concepts to physical models and engineering judgment.
\end{itemize}

\roleentry{NLS Lighting, Carson, CA}{April 2022 -- July 2023}{Mechanical Engineer}{}
\begin{itemize}
  \item Led design-to-release luminaire product improvements across thermal testing, IP testing, FEA static simulation, fabrication, assembly, and installation support.
  \item Produced complete product design packages, including 2D drawings, bills of materials, engineering change orders, installation instructions, and application documentation.
  \item Collaborated with Washington, DC stakeholders to develop a retrofit kit for upgrading more than 75,000 street and alley lights to energy-efficient LEDs as part of a multi-million-dollar Smart Street Lighting project.
\end{itemize}

\roleentry{Hubbell, City of Industry, CA}{May 2021 -- Aug 2021}{Mechanical Engineer Intern}{}
\begin{itemize}
  \item Managed fixture assembly and DLC V5.1 certification planning, reducing test cost and labor while accounting for SKU quantities and sales volume.
  \item Performed impact testing and FEA validation on luminaire assemblies, checking bracket, lens, and LED-PCB designs in Solid Edge.
  \item Supported launch of the Prevent vandal-resistant luminaire series through baseline testing, part sourcing, and assembly coordination with new product design teams.
\end{itemize}

\roleentry{Academy Inc., Montebello, CA}{July 2019 -- Aug 2020}{Mechanical Design Engineer}{}
\begin{itemize}
  \item Implemented parametric product and pricing configurations, reducing design and quoting time.
  \item Engineered shade sails, awnings, and tensioned fabric structures for manufacturability and field installation.
  \item Managed projects from site measurements and material specifications through coordination and scope execution.
\end{itemize}

\roleentry{American Metal Bearing Company, Garden Grove, CA}{May 2019 -- June 2019}{Mechanical Engineer Intern}{}
\begin{itemize}
  \item Optimized thrust, line-shaft, stern-tube, and strut bearing components using SolidWorks/Abaqus FEA, reducing material and labor cost.
  \item Coordinated manufacturing flow for line-shaft spherical bearings from rough casting through inspection, finishing, packaging, and shipment.
  \item Revised bearing-component drawings under Engineering Change Requests while maintaining GD\&T standards.
\end{itemize}

\section*{Technical Skills}
\begin{tabularx}{\textwidth}{@{}>{\bfseries}p{0.18\textwidth} X@{}}
Mechanical design & SolidWorks, Onshape, PTC Creo, FEA review, GD\&T-aware drawings, BOMs, ECOs, DFM, assembly planning, fixture design, prototype iteration.\\
Soft robotics & Pneumatic actuation, compliant mechanisms, continuum robots, modular soft robots, morphing surfaces, airflow routing, pneumatic surface prototypes, simulation-guided robotic prototyping.\\
Software & CUBIT, Sierra/SolidMechanics, ParaView, MATLAB, Simulink Simscape Multibody, Abaqus, Ansys FEA, Python, C/C++, ROS2, LabVIEW, Minitab, MS Office; AI-assisted design and web-tool prototyping.\\
Experiments & Modal testing, vibroacoustic measurement, FFT-based signal analysis, anechoic-chamber testing, optical tracking, pressure/force/displacement characterization, connector/airflow checks, design of experiments.\\
Hardware & 3D printing, TPU/PETG/PLA prototyping, laser cutting, soldering, Arduino, Raspberry Pi, National Instruments DAQ, strain~gages, thermocouples, pressure sensors, force testing (Mecmesin), hand tools.\\
\end{tabularx}

\section*{Relevant Coursework}
Soft Robotics; Advanced Dynamics; Ethics and Communication in Robotics and AI Research; Foundations in Robotics; Advanced Mechanics of Deformable Bodies; Mechatronic System Design; Design of Experiments; Engineering Vibrations; Automatic Feedback Control; Instrumentation and Measurement Systems.

\section*{Campus and Community Involvement}
\datedline{Member, American Society of Mechanical Engineers}{March 2015 -- Present}
\datedline{WPI Video Game Orchestra}{Nov 2024 -- Present}
\datedline{WPI Tennis Club}{Aug 2024 -- Present}
\datedline{WPI Philharmonic Orchestra, 1st Violinist}{Nov 2023 -- Present}
\datedline{CSULB Long Beach Lunabotics, Mechanical Engineer}{Sept 2020 -- May 2022}
\datedline{CSULB Beach Launch Team, Avionics Engineer}{Sept 2020 -- May 2022}
\datedline{CSULB Bob Cole Symphony Orchestra, 2nd Violinist}{Sept 2019 -- March 2020}
\datedline{Rebuilding Together Solano County, Volunteer}{Aug 2015 -- Aug 2016}

\end{document}
